\daytitle{17}{带电粒子在电场中的运动}

\begin{learningbox}{本节导读}
把前两天的电场知识和运动学知识结合起来——带电粒子在电场中会怎样运动？
加速电场中粒子的速度变化，偏转电场中的类平抛运动。示波管的原理也基于此。
\end{learningbox}

\textbf{本节读题路线：} 先判断粒子带正电还是负电，再画电场力方向；加速区优先用动能定理，偏转区把运动分解为水平匀速和竖直匀加速。若题目先加速后偏转，就把两个过程分开列式，再用出口速度衔接。

\section*{一、学习目标}

\begin{enumerate}
  \item 能用动能定理分析加速电场：从静止出发且电场力做正功时，$\lvert q\rvert U = \dfrac12 mv^2$。
  \item 能把偏转电场中的粒子运动分解为水平匀速和竖直匀加速。
  \item 会在图中标出偏转位移、偏转角和离开电场后的运动方向。
  \item 能用“先加速、后偏转”的分段模型说明示波管的基本原理。
  \item 能把“先加速后偏转”综合题分成加速区、偏转区和出场后运动三段，并说明每段使用的规律。
\end{enumerate}
\section*{二、初高中衔接}

初中没有涉及带电粒子运动的内容。本节是电场和运动学的第一次综合——把 Day 02-04 的运动学公式和 Day 15-16 的电场知识结合起来。

\textbf{核心思路：}
\[
\text{电场力 } \vect F=q\vect E
\;\to\;
\text{加速度 } \vect a=\frac{q\vect E}{m}
\]
求大小时使用 $F=|q|E$、$a=|q|E/m$；方向由电荷正负和电场方向判断。

\textbf{科学发展视角：} 带电粒子在电场中的加速和偏转，是电子束、离子注入、粒子探测器和质谱分析的基础。高中阶段先把它看成“电场力 + 运动学”的组合模型，后面进入磁场后，就能理解科学仪器如何用电场和磁场一起控制粒子轨迹。

\section*{三、核心概念}

\begin{conceptbox}{加速电场}

带电粒子在电场中运动时，电场力可能做正功，也可能做负功。本节统一用
$q$ 表示带符号电荷量，用 $|q|$ 表示电荷量大小。电场力做正功时粒子动能增加，做负功时粒子动能减少。

\begin{center}
% TikZ 图：加速电场
% Day 17：带电粒子在电场中的运动
\begin{tikzpicture}[scale=0.9,>=Latex]
  % 加速电极
  \draw[thick] (-2,1) -- (-2,-1);
  \node[left] at (-2,0) {$+$};
  \draw[thick] (2,1) -- (2,-1);
  \node[right] at (2,0) {$-$};
  % 电场线
  \foreach \y in {-0.6,0,0.6} {
    \draw[->,blue] (-1.8,\y) -- (1.8,\y);
  }
  % 电子
  \filldraw[red] (-1.5,0) circle (2pt);
  \draw[->,thick,red] (-1.5,0.4) -- (1.5,0.4) node[above] {$v$};
  \node[below] at (0,-1.3) {$qU = \frac12 mv^2$（初速为零）};
\end{tikzpicture}


\figcaption{加速电场：电场力做功转化为粒子动能}\label{fig:accelerating-field}
\end{center}

\textbf{通式：} 只有电场力做功时，由第16节
$W_{AB}=qU_{AB}$ 和动能定理可得
\[
\boxed{qU_{AB}=\frac12mv_B^2-\frac12mv_A^2}
\]

\textbf{特例（初速为零，电场力做正功）：} 若 $U$ 表示加速电压的大小，则
\[
\boxed{\lvert q\rvert U = \frac12 mv^2}
\]

\textbf{公式从哪里来：} 电场力在加速过程中做的正功为 $W_{\text{电}}=\lvert q\rvert U$，再由动能定理 $W_{\text{电}}=\Delta E_k$，并代入初速度为零的条件，就得到上式。它本质上是 Day 8 的动能定理在电场中的应用。

\begin{itemize}
  \item $U$：加速电压。
  \item $\lvert q\rvert$：粒子电荷量大小。若使用带符号电荷量，应先判断电场力做功正负。
  \item 在只看电场力做功、忽略碰撞和相对论效应的模型中，粒子的末速度 $v = \sqrt{\dfrac{2\lvert q\rvert U}{m}}$ 与加速距离无关。
\end{itemize}

如果粒子有初速度 $v_0$，且电场力仍做正功：
\[
\lvert q\rvert U = \frac12 mv^2 - \frac12 mv_0^2
\]
这里的 $\lvert q\rvert U$ 表示电场力做的正功；若电场力做负功，粒子动能减少，应改用带符号的 $W_{\text{电}}=\Delta E_k$。

\end{conceptbox}

\begin{conceptbox}{偏转电场（类平抛）}

带电粒子以初速度 $v_0$ 垂直射入匀强电场，做类平抛运动。以下统一取 $x$ 轴正向为初速度方向，$y$ 轴正向为电场方向；$E$ 和 $U$ 表示场强、电压的大小。这样 $q$ 保留正负号，$y$、$a_y$ 和偏转角都可带正负号。

\begin{center}
% TikZ 图：偏转电场（类平抛）
% Day 17：带电粒子在电场中的运动
\begin{tikzpicture}[scale=0.9,>=Latex]
  % 平行板
  \draw[thick] (0,1.5) -- (3,1.5);
  \node[above] at (1.5,1.5) {$+$};
  \draw[thick] (0,-1.5) -- (3,-1.5);
  \node[below] at (1.5,-1.5) {$-$};
  % 电场线
  \foreach \x in {0.5,1.5,2.5} {
    \draw[->,blue] (\x,1.3) -- (\x,-1.3);
  }
  % 入射正电粒子
  \filldraw[red] (-1,0) circle (2pt);
  \node[below] at (-1,0) {$+q$};
  \draw[->,thick,red] (-1,0) -- (0.5,0) node[above] {$v_0$};
  % 轨迹（抛物线）
  \draw[thick,red,dashed,domain=0:3.45,samples=80]
    plot (\x,{-0.12*\x*\x});
  \filldraw[red] (3.45,-1.43) circle (2pt);
  % 偏转量
  \draw[dashed,gray] (3.45,0) -- (3.45,-1.43);
  \draw[<->] (3.75,0) -- (3.75,-1.43);
  \node[right] at (3.75,-0.72) {$y$};
  \draw[->,red] (3.45,-1.43) -- (4.25,-2.10) node[right] {$v$};
  \node[red,below] at (1.65,-0.35) {抛物线};
\end{tikzpicture}


\figcaption{偏转电场：粒子轨迹为抛物线}\label{fig:deflection-field}
\end{center}

\textbf{分析方法：} 正交分解
\begin{itemize}
  \item $x$ 方向（匀速）：$x = v_0 t$。
  \item $y$ 方向（匀加速）：$a_y = \dfrac{qE}{m} = \dfrac{qU}{md}$，$y = \dfrac12 a_y t^2$。$q>0$ 时粒子沿电场方向偏转，$q<0$ 时反向偏转；若只求大小，则用 $\lvert a_y\rvert=\dfrac{|q|E}{m}$。
\end{itemize}

\textbf{为什么轨迹是抛物线：} 由 $x=v_0t$ 得
$t=x/v_0$，代入 $y=\frac12a_yt^2$：
\[
\boxed{y=\frac{qE}{2mv_0^2}x^2}
\]
这是关于 $x$ 的二次函数，因此有场区内的轨迹为抛物线；$q$ 的正负决定抛物线开口方向。

\textbf{偏转位移（极板长 $L$，粒子从极板中心射入，忽略重力和边缘场）：}
\[
y = \frac12 \cdot \frac{qU}{md} \cdot \left(\frac{L}{v_0}\right)^2
\]

\textbf{偏转角 $\theta$（出射方向与水平方向的夹角）：}
\[
\tan\theta = \frac{v_y}{v_0} = \frac{qUL}{md v_0^2}
\]

\textbf{不撞板条件：} 粒子从两板中央射入时，离开极板前应满足 $\lvert y\rvert<\dfrac{d}{2}$；若不满足，粒子会先撞到极板，不能再套用“完整穿过极板”的公式。

\end{conceptbox}

\begin{samepage}
\section*{四、公式与适用条件}

\textbf{方法使用前检查：}
\nopagebreak[4]
\begin{enumerate}
  \item 先判断粒子电性，确定电场力方向，再决定加速度方向。
  \item 加速区优先用动能定理；偏转区把运动分解为水平方向和竖直方向。
  \item 先规定坐标正方向。以下公式取 $y$ 轴正向为电场方向，$E$、$U$ 取大小，$q$ 保留正负号。
  \item 若只受电场力，竖直方向可能是匀加速运动，水平方向可能仍是匀速运动；还要检查粒子会不会先撞到极板。
  \item 出场后进入无场区时，速度方向保持出场瞬间方向，轨迹变为直线。
\end{enumerate}
\end{samepage}

\begin{center}
\small
\begin{tabular}{>{\centering\arraybackslash}p{0.31\textwidth}p{0.20\textwidth}p{0.36\textwidth}}
\toprule
公式 & 物理意义 & 适用条件 \\
\midrule
$qU_{AB}=\Delta E_k$ & 电场加速通式 & 只有电场力做功；$q$ 和 $U_{AB}$ 均按定义带符号 \\
$\lvert q\rvert U = \frac12 mv^2$ & 特例：从静止加速 & 电场力做正功，$U$ 取电压大小 \\
$a_y = \dfrac{qE}{m} = \dfrac{qU}{md}$ & 电场中加速度 & 匀强电场；$y$ 正向取电场方向，$q$ 可带符号 \\
$y = \frac12 \dfrac{qU}{md}\left(\dfrac{L}{v_0}\right)^2$ & 偏转位移 & 垂直射入且未撞板 \\
$\tan\theta = \dfrac{qUL}{mdv_0^2}$ & 有向偏转角 & 同上 \\
\bottomrule
\end{tabular}
\end{center}

\begin{learningbox}{带电粒子运动的分段分析清单}
带电粒子题先按区域切开，每一段只抓一个主工具：
\begin{center}
\begin{tabular}{p{0.18\textwidth}p{0.30\textwidth}p{0.40\textwidth}}
\toprule
区域 & 主工具 & 关键判断 \\
\midrule
加速区 & 动能定理 $qU_{AB}=\Delta E_k$ & 先判断电场力是否做正功；从静止加速时才可直接写 $\lvert q\rvert U=\frac12mv^2$。 \\
偏转区 & 类平抛分解 & 沿入射方向通常匀速，沿电场力方向匀加速；通过时间由 $t=L/v_0$ 决定。 \\
无场区 & 匀速直线运动 & 离开极板后不再受电场力，速度大小和方向保持出场瞬间状态，沿出射速度方向运动到屏幕。 \\
\bottomrule
\end{tabular}
\end{center}
\textbf{读题提醒：} 先由电荷正负确定电场力方向，再把上一段的末速度作为下一段的初速度。不要把加速电压 $U_1$ 和偏转电压 $U_2$ 混在同一个步骤里。
\end{learningbox}

\section*{五、典型模型}

\subsection*{模型 1：先加速后偏转}

粒子先经过加速电压 $U_1$ 加速，再进入偏转电压 $U_2$ 的偏转电场。

\begin{center}
% TikZ 图：先加速后偏转模型
% Day 17：带电粒子在电场中的运动
\begin{tikzpicture}[scale=0.78,>=Latex]
  % 加速区
  \draw[thick] (-3.2,1.0) -- (-3.2,-1.0);
  \draw[thick] (-1.6,1.0) -- (-1.6,-1.0);
  \node[above] at (-3.2,1.04) {$+$};
  \node[above] at (-1.6,1.04) {$-$};
  \foreach \y in {-0.45,0,0.45} {
    \draw[->,blue] (-3.0,\y) -- (-1.8,\y);
  }
  \draw[<->] (-3.2,1.32) -- (-1.6,1.32);
  \node[above,blue,fill=white,inner sep=1pt] at (-2.4,1.32) {$U_1$};
  \node[below left,red,fill=white,inner sep=1pt] at (-3.6,0) {$+q$};
  \draw[red,thick,dashed] (-3.6,0) -- (-1.62,0);
  \draw[-{Stealth[length=5pt,width=6pt]},red,line width=1.05pt] (-1.35,0) -- (-0.35,0);
  \node[red,fill=white,inner sep=1pt] at (-0.82,0.28) {$v_0$};

  % 偏转区
  \draw[thick] (-0.2,1.2) -- (3.0,1.2);
  \draw[thick] (-0.2,-1.2) -- (3.0,-1.2);
  \node[above] at (-0.35,1.2) {$+$};
  \node[below] at (-0.35,-1.2) {$-$};
  \foreach \x in {0.35,1.05,1.75,2.45} {
    \draw[->,blue] (\x,1.0) -- (\x,-1.0);
  }
  \draw[<->] (-0.2,1.55) -- (3.0,1.55);
  \node[above,fill=white,inner sep=1pt] at (1.4,1.55) {$L$};
  \draw[<->] (3.35,1.2) -- (3.35,-1.2);
  \node[right,fill=white,inner sep=1pt] at (3.35,0) {$d$};
  \node[blue,fill=white,inner sep=1pt] at (0.15,-1.48) {$U_2$};

  \draw[red,thick,dashed,domain=0.25:3.3,samples=80]
    plot (\x,{-0.10*(\x+0.2)*(\x+0.2)});
  \draw[dashed,gray] (3.3,0) -- (3.3,-1.23);
  \draw[<->] (4.0,0) -- (4.0,-1.23);
  \node[right,fill=white,inner sep=1pt] at (4.0,-0.62) {$y$};

  \node[font=\small,fill=white,inner sep=1pt] at (-2.4,-1.48) {加速区};
  \node[font=\small,fill=white,inner sep=1pt] at (2.15,-1.82) {偏转区};
\end{tikzpicture}


\figcaption{先加速后偏转模型：$U_1$ 决定入射速度，$U_2$ 决定偏转大小}\label{fig:accelerate-deflect-model}
\end{center}

加速：$\lvert q\rvert U_1 = \frac12 mv_0^2 \;\Rightarrow\; v_0^2 = \frac{2\lvert q\rvert U_1}{m}$

偏转量大小：
\[
\boxed{|y| = \frac{U_2 L^2}{4U_1 d}}
\]
相应的偏转角大小满足
\[
\boxed{|\tan\theta|=\frac{U_2L}{2U_1d}}.
\]
在粒子从静止经同一加速电压加速、进入同一理想偏转电场并忽略重力和相对论效应时，偏转量和偏转角的
\textbf{大小}与 $|q|$、$m$ 无关；偏转方向仍由电荷正负和偏转板极性决定。

\textbf{分段分析清单：} 加速区只负责求入射偏转区的速度 $v_0$；偏转区用这个 $v_0$ 求通过时间 $t=L/v_0$；竖直方向再用 $a_y=qU_2/(md)$ 求偏转位移。不要一开始就把两个电压混在同一个公式里。

\subsection*{模型 2：示波管}

示波管可以看成“先加速后两次偏转”的装置。电子枪发射电子，电子先经加速电场获得较大的水平速度，再依次经过两对偏转板，最后打在荧光屏上形成亮点。

\begin{center}
% TikZ 图：示波管模型
% Day 17：带电粒子在电场中的运动
\begin{tikzpicture}[scale=0.76,>=Latex]
  % 电子枪与加速区
  \draw[thick] (-4.45,0.55) -- (-3.85,0) -- (-4.45,-0.55);
  \node[left] at (-4.45,0) {电子枪};
  \draw[thick] (-3.0,0.9) -- (-3.0,-0.9);
  \draw[thick] (-2.35,0.9) -- (-2.35,-0.9);
  \foreach \y in {-0.35,0,0.35} {
    \draw[->,blue] (-2.9,\y) -- (-2.45,\y);
  }
  \node[below,font=\small,fill=white,inner sep=1pt] at (-2.68,-1.05) {加速区};

  % y 偏转板
  \draw[thick] (-1.35,0.8) -- (0.55,0.8);
  \draw[thick] (-1.35,-0.8) -- (0.55,-0.8);
  \node[above,font=\small,fill=white,inner sep=1pt] at (-0.4,0.9) {$Y$偏转板};
  \foreach \x in {-1.0,-0.35,0.3} {
    \draw[->,blue!65] (\x,0.65) -- (\x,-0.65);
  }

  % x 偏转板（侧视中用窄板表示）
  \draw[thick] (1.05,0.65) -- (1.05,-0.65);
  \draw[thick] (1.85,0.65) -- (1.85,-0.65);
  \node[below,font=\small,fill=white,inner sep=1pt] at (1.45,-1.02) {$X$偏转板};

  % 屏幕
  \draw[thick] (3.7,1.45) -- (3.7,-1.45);
  \node[right,fill=white,inner sep=1pt] at (3.7,1.2) {荧光屏};
  \draw[red,thick,dashed] (-3.85,0) -- (-1.35,0)
    .. controls (-0.55,-0.05) and (0.25,-0.35) .. (0.7,-0.45)
    -- (1.9,-0.45)
    .. controls (2.45,-0.40) and (3.05,-0.20) .. (3.7,0.25);
  \filldraw[red] (3.7,0.25) circle (2.3pt);
  \node[right,red,fill=white,inner sep=1pt] at (3.9,0.25) {亮点};

\end{tikzpicture}


\figcaption{示波管模型：两组偏转板相互垂直，分别用侧视图和俯视图表示}\label{fig:oscilloscope-model}
\end{center}

\begin{itemize}
  \item \textbf{加速电场：} 决定电子进入偏转系统前的速度 $v_0$。电子带负电，受力方向与电场方向相反。
  \item \textbf{$Y$ 偏转板：} 控制亮点的上下位置，通常接入待观测的信号电压。
  \item \textbf{$X$ 偏转板：} 控制亮点的左右位置，常接扫描电压。
  \item \textbf{荧光屏：} 电子打到屏上形成亮点，亮点随时间移动就形成图像。
\end{itemize}

\subsection*{模型 3：离开偏转电场后的无场区运动}

粒子在偏转电场内做类平抛运动；离开极板后若进入无场区，就沿出射速度方向做匀速直线运动。因此屏幕上的总偏移通常由两部分组成：
\[
Y = y_1 + y_2
\]
其中 $y_1$ 是极板内偏转，$y_2$ 是离开极板后在无场区沿出射速度方向匀速运动产生的偏移。

设偏转板长为 $L$，极板出口到屏的水平距离为 $s$。在极板内，
\[
y_1=\frac{qEL^2}{2mv_0^2},
\qquad
v_y=\frac{qEL}{mv_0}.
\]
离开极板后运动时间为 $s/v_0$，故
\[
y_2=v_y\frac{s}{v_0}
=\frac{qELs}{mv_0^2}.
\]
于是屏上的有向总偏移为
\[
\boxed{Y=y_1+y_2
=\frac{qEL}{mv_0^2}\left(s+\frac L2\right)}.
\]
若只求偏移大小，应取绝对值；方向仍由 $q\vect E$ 判断。

\begin{center}
% TikZ 图：离开偏转电场后的无场区运动
% Day 17：带电粒子在电场中的运动
\begin{tikzpicture}[scale=0.74,>=Latex]
  \draw[thick] (0,1.15) -- (3.0,1.15);
  \draw[thick] (0,-1.15) -- (3.0,-1.15);
  \node[above] at (-0.15,1.15) {$+$};
  \node[below] at (-0.15,-1.15) {$-$};
  \foreach \x in {0.45,1.15,1.85,2.55} {
    \draw[->,blue] (\x,0.95) -- (\x,-0.95);
  }
  \draw[red,thick,dashed,domain=0:3,samples=70] plot (\x,{-0.075*\x*\x});
  \coordinate (P) at (3,-0.675);
  \filldraw[red] (P) circle (1.8pt);
  \coordinate (S) at (5.8,-1.935);
  \draw[red,thick,dashed] (4.08,-1.16) -- (S);
  \draw[-{Stealth[length=5pt,width=6pt]},red,line width=1.05pt,shorten <=2pt] (P) -- (4.08,-1.16);
  \node[above right,red,fill=white,inner sep=1pt] at (4.08,-1.02) {$v$};
  \draw[dashed,gray] (3,0) -- (P);
  \draw[dashed,gray] (5.8,0) -- (S);
  \draw[<->] (3.65,0) -- (3.65,-0.675);
  \node[right,fill=white,inner sep=1pt] at (3.72,-0.34) {$y_1$};
  \draw[dashed,gray] (3,-0.675) -- (5.8,-0.675);
  \draw[<->] (5.35,-0.675) -- (5.35,-1.935);
  \node[left,fill=white,inner sep=1pt] at (5.35,-1.30) {$y_2$};
  \draw[<->] (6.25,0) -- (6.25,-1.95);
  \node[right,fill=white,inner sep=1pt] at (6.25,-0.98) {$Y$};
  \draw[<->] (3,-2.08) -- (5.8,-2.08);
  \node[below,fill=white,inner sep=1pt] at (4.4,-2.08) {$s$};
  \draw[thick] (5.8,1.1) -- (5.8,-2.2);
  \node[right,fill=white,inner sep=1pt] at (5.8,1.0) {屏};
  \node[font=\small,fill=white,inner sep=1pt] at (1.5,1.55) {有场区};
  \node[font=\small,fill=white,inner sep=1pt] at (4.4,0.55) {无场区};
\end{tikzpicture}


\figcaption{离开偏转电场后的无场区：沿出射速度方向匀速直线运动}\label{fig:post-deflection-drift}
\end{center}

\newpage
\section*{六、例题精讲}

\begin{examplebox}{加速电场}

\textbf{题意分析：} 一个电子从静止开始被 $U = 500\ \mathrm{V}$ 的电压加速。已知 $e = 1.6 \times 10^{-19}\ \mathrm{C}$，$m_e = 9.1 \times 10^{-31}\ \mathrm{kg}$。求电子加速后的速度。

\textbf{物理模型：} 加速电场，$\lvert q\rvert U = \frac12 mv^2$。

\textbf{解题过程：}

\[
\begin{aligned}
v &= \sqrt{\frac{2eU}{m}}
= \sqrt{\frac{2 \times 1.6 \times 10^{-19} \times 500}{9.1 \times 10^{-31}}} \\
&= \sqrt{1.76 \times 10^{14}}
\approx 1.33 \times 10^7\ \mathrm{m/s}
\end{aligned}
\]

\textbf{方法总结：} 电子的质量很小，$500\ \mathrm{V}$ 的电压就能让它加速到 $10^7\ \mathrm{m/s}$ 量级。但即使如此，远小于光速 $3 \times 10^8\ \mathrm{m/s}$，不需用相对论修正。

\end{examplebox}

\begin{examplebox}{偏转电场}

\textbf{题意分析：} 电子以 $v_0 = 3 \times 10^7\ \mathrm{m/s}$ 水平射入两平行板间。极板长 $L = 0.06\ \mathrm{m}$，间距 $d = 0.02\ \mathrm{m}$，电压 $U = 200\ \mathrm{V}$。求：
\begin{enumerate}
  \item 电子在电场中的运动时间。
  \item 电子的偏转位移。
  \item 出射时的偏转角。
\end{enumerate}

\textbf{物理模型：} 类平抛运动。

\textbf{解题过程：}

\begin{enumerate}
  \item $t = \frac{L}{v_0} = \frac{0.06}{3 \times 10^7} = 2 \times 10^{-9}\ \mathrm{s}$

  \item 由图~\ref{fig:deflection-field} 所示极板极性，电场方向向下。取向下为 $y$ 轴正向，电子 $q=-e$，故 $a_y=-\dfrac{eU}{md}$，实际向上偏转；其加速度大小为
  \[
  \lvert a_y\rvert = \frac{eU}{md} = \frac{1.6 \times 10^{-19} \times 200}{9.1 \times 10^{-31} \times 0.02} \approx 1.76 \times 10^{15}\ \mathrm{m/s^2}.
  \]

  因而 $y=-\dfrac12\lvert a_y\rvert t^2=-3.52\ \mathrm{mm}$，即向上偏转 $3.52\ \mathrm{mm}$。又 $\lvert y\rvert=3.52\ \mathrm{mm}<d/2=10\ \mathrm{mm}$，粒子不会撞板。

  \item $\tan\theta = \dfrac{v_y}{v_0} = \dfrac{a_yt}{v_0}\approx -0.117$

  $\theta \approx -6.7^\circ$；负号表示相对初速度方向向上偏转，偏转角大小为 $6.7^\circ$。
\end{enumerate}

\textbf{方法总结：} 偏转电场的关键是把运动分解为匀速和匀加速两个方向分别处理，数学处理与平抛运动相同，但加速度方向要由电荷正负和电场方向判断。

\end{examplebox}

\begin{examplebox}{先加速后偏转}

\textbf{题意分析：} 电子从静止经 $U_1 = 1000\ \mathrm{V}$ 加速，然后垂直进入 $U_2 = 50\ \mathrm{V}$ 的偏转电场。极板长 $L = 0.05\ \mathrm{m}$，间距 $d = 0.01\ \mathrm{m}$。电子从两板中央射入，求电子离开极板时的偏转量，并判断是否撞板。

\textbf{物理模型：} 先加速后偏转。

\begin{center}
% TikZ 图：先加速后偏转例题
% Day 17：带电粒子在电场中的运动
\begin{tikzpicture}[scale=0.82,>=Latex]
  \draw[thick] (-3.0,0.9) -- (-3.0,0.16);
  \draw[thick] (-3.0,-0.16) -- (-3.0,-0.9);
  \node[above] at (-3.0,0.9) {$-$};
  \draw[thick] (-1.55,0.9) -- (-1.55,0.16);
  \draw[thick] (-1.55,-0.16) -- (-1.55,-0.9);
  \node[above] at (-1.55,0.9) {$+$};
  \foreach \y in {-0.4,0,0.4} {
    \draw[->,blue] (-1.75,\y) -- (-2.8,\y);
  }
  \node[below,fill=white,inner sep=1pt] at (-2.25,-1.05) {$U_1=1000\,\mathrm{V}$};
  \filldraw[red] (-3.35,0) circle (2pt);
  \node[below left,red,fill=white,inner sep=1pt] at (-3.35,0) {$e^-$};
  \draw[red,thick,dashed] (-3.20,0) -- (-1.57,0);
  \draw[-{Stealth[length=5pt,width=6pt]},red,line width=1.05pt] (-1.32,0) -- (-0.45,0);
  \node[red,fill=white,inner sep=1pt] at (-0.86,0.28) {$v_0$};

  \draw[thick] (-0.1,1.15) -- (2.9,1.15);
  \draw[thick] (-0.1,-1.15) -- (2.9,-1.15);
  \node[above] at (1.4,1.15) {$-$};
  \node[below] at (1.4,-1.15) {$+$};
  \foreach \x in {0.45,1.15,1.85,2.55} {
    \draw[->,blue] (\x,-0.95) -- (\x,0.95);
  }
  \node[below,fill=white,inner sep=1pt] at (0.85,-1.45) {$U_2=50\,\mathrm{V}$};
  \draw[<->] (-0.1,1.45) -- (2.9,1.45);
  \node[above] at (1.4,1.45) {$L=0.05\,\mathrm{m}$};
  \draw[<->] (3.18,1.15) -- (3.18,-1.15);
  \node[right,fill=white,inner sep=1pt] at (3.18,0.35) {$d=0.01\,\mathrm{m}$};
  \draw[red,thick,dashed,domain=-0.1:2.9,samples=80]
    plot (\x,{-0.08*(\x+0.1)*(\x+0.1)});
  \filldraw[red] (2.9,-0.72) circle (1.8pt);
  \draw[dashed,gray] (2.9,0) -- (2.9,-0.72);
  \draw[<->] (3.78,0) -- (3.78,-0.72);
  \node[right,fill=white,inner sep=1pt] at (3.78,-0.36) {$|y|$};
\end{tikzpicture}


\figcaption{例题示意：先由 $U_1$ 加速，再由 $U_2$ 偏转}\label{fig:accelerate-deflect-example}
\end{center}

\textbf{解题过程：}

加速后速度 $v_0^2 = \dfrac{2eU_1}{m}$

偏转位移：
\[
|y| = \frac12 \frac{eU_2}{md} \left(\frac{L}{v_0}\right)^2
= \frac12 \frac{eU_2}{md} \cdot \frac{L^2 m}{2eU_1}
= \frac{U_2 L^2}{4U_1 d}
\]

代入：
\[
|y| = \frac{50 \times (0.05)^2}{4 \times 1000 \times 0.01}
=0.003125\ \mathrm{m}=3.125\ \mathrm{mm}.
\]
由于 $|y|=3.125\ \mathrm{mm}<d/2=5.00\ \mathrm{mm}$，电子能从极板间飞出。图示电场方向向上，电子带负电，因此实际向下偏转。

\textbf{方法总结：} 先加速后偏转题不要直接背最终式。先分成两段：加速区用 $\lvert q\rvert U_1=\frac12mv_0^2$ 求入射速度；偏转区用 $x=v_0t$、$y=\frac12a_yt^2$。只有在“从静止加速、理想偏转电场、忽略重力和相对论效应”这些条件同时满足时，偏转量大小才可化简为 $|y| = \dfrac{U_2 L^2}{4U_1 d}$。最后仍须判断方向并检查 $|y|<d/2$。

\end{examplebox}

\section*{七、变式训练}

\begin{enumerate}
  \item 质子（$m_p = 1.67 \times 10^{-27}\ \mathrm{kg}$，$q = e$）从静止经 $U = 1000\ \mathrm{V}$ 加速，求末速度。

  \item 电子从两板中线以 $v_0 = 2 \times 10^7\ \mathrm{m/s}$ 射入图示偏转电场，$L = 0.04\ \mathrm{m}$，$d = 0.01\ \mathrm{m}$，$U = 150\ \mathrm{V}$。先判断电子能否从极板右端飞出；若不能，求碰撞点到入口的水平距离。
  \begin{center}
  \begin{tikzpicture}[scale=0.78,>=Stealth]
  \draw[thick] (0,1.45) -- (4.4,1.45);
  \draw[thick] (0,0) -- (4.4,0);
  \node[left] at (0,1.45) {$-$};
  \node[left] at (0,0) {$+$};
  \foreach \x in {0.7,1.5,2.3,3.1,3.9}
    \draw[->,blue] (\x,0.25) -- (\x,1.25);
  \draw[->,red,thick] (-0.55,0.72) -- (0.35,0.72) node[midway,above] {$v_0$};
  \node[red,left] at (-0.55,0.72) {$e^-$};
  \draw[red,thick,dashed,domain=0.35:4.25,samples=80] plot (\x,{0.72-0.035*(\x-0.35)^2});
  \draw[<->] (0,-0.35) -- (4.4,-0.35) node[midway,below] {$L$};
  \draw[<->] (4.7,0) -- (4.7,1.45) node[midway,right] {$d$};
\end{tikzpicture}

  \end{center}

  \item 质子和 $\alpha$ 粒子（$m_\alpha = 4m_p$，$q_\alpha = 2e$）经过同一加速电压 $U_1$ 加速后进入同一偏转电场，比较它们的偏转位移和偏转角。
  \begin{center}
  % TikZ 图：先加速后偏转模型
% Day 17：带电粒子在电场中的运动
\begin{tikzpicture}[scale=0.78,>=Latex]
  % 加速区
  \draw[thick] (-3.2,1.0) -- (-3.2,-1.0);
  \draw[thick] (-1.6,1.0) -- (-1.6,-1.0);
  \node[above] at (-3.2,1.04) {$+$};
  \node[above] at (-1.6,1.04) {$-$};
  \foreach \y in {-0.45,0,0.45} {
    \draw[->,blue] (-3.0,\y) -- (-1.8,\y);
  }
  \draw[<->] (-3.2,1.32) -- (-1.6,1.32);
  \node[above,blue,fill=white,inner sep=1pt] at (-2.4,1.32) {$U_1$};
  \node[below left,red,fill=white,inner sep=1pt] at (-3.6,0) {$+q$};
  \draw[red,thick,dashed] (-3.6,0) -- (-1.62,0);
  \draw[-{Stealth[length=5pt,width=6pt]},red,line width=1.05pt] (-1.35,0) -- (-0.35,0);
  \node[red,fill=white,inner sep=1pt] at (-0.82,0.28) {$v_0$};

  % 偏转区
  \draw[thick] (-0.2,1.2) -- (3.0,1.2);
  \draw[thick] (-0.2,-1.2) -- (3.0,-1.2);
  \node[above] at (-0.35,1.2) {$+$};
  \node[below] at (-0.35,-1.2) {$-$};
  \foreach \x in {0.35,1.05,1.75,2.45} {
    \draw[->,blue] (\x,1.0) -- (\x,-1.0);
  }
  \draw[<->] (-0.2,1.55) -- (3.0,1.55);
  \node[above,fill=white,inner sep=1pt] at (1.4,1.55) {$L$};
  \draw[<->] (3.35,1.2) -- (3.35,-1.2);
  \node[right,fill=white,inner sep=1pt] at (3.35,0) {$d$};
  \node[blue,fill=white,inner sep=1pt] at (0.15,-1.48) {$U_2$};

  \draw[red,thick,dashed,domain=0.25:3.3,samples=80]
    plot (\x,{-0.10*(\x+0.2)*(\x+0.2)});
  \draw[dashed,gray] (3.3,0) -- (3.3,-1.23);
  \draw[<->] (4.0,0) -- (4.0,-1.23);
  \node[right,fill=white,inner sep=1pt] at (4.0,-0.62) {$y$};

  \node[font=\small,fill=white,inner sep=1pt] at (-2.4,-1.48) {加速区};
  \node[font=\small,fill=white,inner sep=1pt] at (2.15,-1.82) {偏转区};
\end{tikzpicture}

\end{center}
\end{enumerate}

\begin{learningbox}{变式训练参考结果}
\begin{enumerate}[itemsep=0.45em,topsep=0.3em,parsep=0pt]
  \item 由 $eU=\frac12m_pv^2$，得
  $v\approx4.38\times10^5\ \mathrm{m/s}$。
  \item 若先假设电子能到达极板末端，所需时间
  $t=L/v_0=2.0\times10^{-9}\ \mathrm s$，形式计算得到
  \[
  y=\frac12\frac{eU}{m_ed}t^2\approx5.27\times10^{-3}\ \mathrm m,
  \]
  但 $|y|>d/2=5.00\ \mathrm{mm}$，电子会在到达极板末端前撞上下极板。由
  \[
  \frac d2=\frac12\frac{eU}{m_ed}t_{\text{碰}}^2
  \]
  得 $t_{\text{碰}}\approx1.95\times10^{-9}\ \mathrm s$，故碰撞点到入口的水平距离
  $x_{\text{碰}}=v_0t_{\text{碰}}\approx3.89\times10^{-2}\ \mathrm m$。
  \item 经同一电压加速后，$v_0^2=2qU_1/m$。代入偏转公式可得
  $|y|=EL^2/(4U_1)$、$|\tan\theta|=EL/(2U_1)$，均与粒子的质量和电荷量大小无关。质子和 $\alpha$ 粒子均带正电，因此偏转量、偏转角的大小和方向都相同。
\end{enumerate}
\end{learningbox}

\section*{八、竞赛入门补充}

\begin{competitionbox}
\textbf{带电粒子在交变电场中的运动}

如果偏转电场是交变的（如示波器中的锯齿波或正弦波电压），粒子的偏转位移随时间变化，在荧光屏上形成图像。

\begin{center}
% TikZ 图：交变偏转电场与示波图像
% Day 17：带电粒子在电场中的运动
\begin{tikzpicture}[scale=0.72,>=Latex]
  % 左侧：输入电压随时间
  \draw[->] (-3.8,0) -- (-0.4,0) node[right] {$t$};
  \draw[->] (-3.6,-1.1) -- (-3.6,1.1) node[above] {$U_y$};
  \draw[red,thick,domain=-3.5:-0.55,samples=90]
    plot (\x,{0.72*sin(240*(\x+3.5))});
  \node[red] at (-2.0,1.05) {正弦信号};

  \draw[->] (-3.8,-2.3) -- (-0.4,-2.3) node[right] {$t$};
  \draw[->] (-3.6,-3.1) -- (-3.6,-1.55) node[above] {$U_x$};
  \draw[blue,thick] (-3.5,-3.0) -- (-2.55,-1.72) -- (-2.55,-3.0)
    -- (-1.6,-1.72) -- (-1.6,-3.0) -- (-0.65,-1.72);
  \node[blue] at (-2.0,-1.45) {扫描电压};

  % 右侧：屏幕上的轨迹
  \draw[rounded corners=2pt,thick] (0.7,-2.6) rectangle (4.3,1.1);
  \draw[->,gray] (0.95,-0.75) -- (3.78,-0.75);
  \node[below right,gray,fill=white,inner sep=1pt] at (3.72,-0.82) {$x$};
  \draw[->,gray] (2.5,-2.35) -- (2.5,0.62);
  \node[left,gray,fill=white,inner sep=1pt] at (2.43,0.58) {$y$};
  \draw[red,thick,domain=0.95:4.05,samples=120]
    plot (\x,{-0.75 + 0.75*sin(210*(\x-0.95))});
  \node at (2.5,1.35) {荧光屏上形成波形};
  \node[font=\footnotesize,align=center] at (2.5,-3.05)
    {$x$ 随时间匀速扫描，$y$ 随信号电压上下偏转};
\end{tikzpicture}


\figcaption{交变偏转电场：扫描电压把时间展开，信号电压决定上下偏转}\label{fig:alternating-field-trace}
\end{center}

\textbf{关键思路：} 粒子穿过偏转电场的时间 $t = L/v_0$ 非常短（通常 $< 10^{-8}\ \mathrm{s}$），而交变电压的周期 $T$ 往往长得多，所以单个粒子通过极板时，可以近似认为电压不变。不同时间发出的电子对应不同瞬时电压，打在屏幕上的位置不同，连续起来就形成图像。

\textbf{例题：} 示波器中 $x$ 偏转板加锯齿波电压（扫描电压），$y$ 偏转板加正弦信号电压。电子打在屏上的轨迹是什么？

\textbf{解：} $y$ 方向偏转与正弦电压成正比，亮点随信号电压上下运动；$x$ 方向偏转与锯齿波扫描电压成正比，亮点随时间近似匀速向右运动。因此屏上把“随时间变化的正弦电压”展开成一条正弦曲线。

\textbf{进一步理解：} 如果扫描电压周期与信号周期匹配，每次扫描画出的曲线重合，屏幕上看到稳定波形；如果二者不同步，波形会左右移动或变得不稳定。

\textbf{回扫说明：} 锯齿波突然回落时，亮点会快速返回屏幕左侧。实际仪器通常在回扫阶段使电子束暂时熄灭，避免屏幕上出现干扰回扫线。

\textbf{竞赛提示：} 处理交变电场问题时，先判断粒子经过电场的时间与交变周期的关系。如果 $t \ll T$，可以对单个粒子近似为恒定电场；再把不同时刻的粒子落点按时间顺序连起来。
\end{competitionbox}

\section*{九、易错点}

\begin{mistakebox}
\begin{enumerate}
  \item \textbf{把电场力加速度和重力加速度混淆：} 先判断是否忽略重力。只受电场力时加速度大小 $a = \dfrac{|q|E}{m}$，不是 $g$；方向由电荷正负和电场方向决定。

  \item \textbf{偏转电场中忘记 $x$ 方向是匀速运动：} 先分解到水平和竖直方向。$x$ 方向不受力时是匀速直线运动，不能写成匀加速。

  \item \textbf{混淆加速电压和偏转电压：} 先把过程分成加速区和偏转区。加速 $U_1$ 决定 $v_0$，偏转 $U_2$ 决定竖直加速度和偏移。

  \item \textbf{有初速度时仍套 $\lvert q\rvert U=\frac12mv^2$：} 这个式子只适用于从静止开始、且电场力做正功的情形。若粒子已有初速度，应写成 $W_{\text{电}}=\Delta E_k=\frac12mv^2-\frac12mv_0^2$，并先判断电场力做功正负。

  \item \textbf{忽略重力的条件：} 先看研究对象尺度。电子、质子、离子在电场中通常忽略重力；带电液滴、带电小球则要把重力列入受力分析。

  \item \textbf{偏转公式中的 $L$ 是极板长度：} 先标出极板长度和屏距。极板内偏转用 $L$，出板后的漂移距离要另算。
\end{enumerate}
\end{mistakebox}

\section*{十、今日小结}

\begin{enumerate}
  \item 只有电场力做功时，加速电场的通式是 $qU_{AB}=\Delta E_k$；从静止开始且电场力做正功时，才化为 $\lvert q\rvert U = \frac12 mv^2$。
  \item 偏转电场：类平抛运动，$x$ 方向匀速 $x = v_0 t$，$y$ 方向匀加速 $a_y = \dfrac{qE}{m}$；先规定 $y$ 正方向，再处理正负号。
  \item 偏转位移 $y = \dfrac12 \dfrac{qU}{md}\left(\dfrac{L}{v_0}\right)^2$，与 $q$、$m$ 有关；使用前检查 $\lvert y\rvert<d/2$。
  \item 先加速后偏转：在从静止加速、同一偏转电场且忽略重力和相对论效应时，
  $|y| = \dfrac{U_2 L^2}{4 U_1 d}$，其大小与 $|q|$、$m$ 无关，方向仍由电荷正负和极板极性决定。
  \item 微观粒子忽略重力，宏观带电物体要考虑重力。
\end{enumerate}

\begin{center}
% TikZ 图：带电粒子运动的三段模型总览
\begin{tikzpicture}[scale=0.82,>=Latex,font=\small]
  \node[draw,minimum width=3.0cm,minimum height=1.35cm,align=center] (a) at (-4.4,0)
    {\textbf{加速区}\\[2pt]
     $qU_{AB}=\Delta E_k$\\
     求入射速度 $v_0$};
  \node[draw,minimum width=3.2cm,minimum height=1.35cm,align=center] (b) at (0,0)
    {\textbf{偏转区}\\[2pt]
     $x=v_0t,\quad y=\frac12a_yt^2$\\
     检查 $|y|<d/2$};
  \node[draw,minimum width=3.0cm,minimum height=1.35cm,align=center] (c) at (4.4,0)
    {\textbf{无场区}\\[2pt]
     $\vect v=\text{常量}$\\
     沿出口切线直线运动};

  \draw[->,thick] (a.east) -- (b.west);
  \draw[->,thick] (b.east) -- (c.west);
\end{tikzpicture}

\end{center}

\textbf{一分钟复盘：} 画一个“先加速后偏转”的带电粒子路径，按三段标注：加速区写能量方程，偏转区写 $x$、$y$ 两方向运动规律，离开电场后画速度方向并写匀速直线运动。
